\scanpage{580}

\mspage{489}

a) Soit $x\in\mathrm{Ob}\,M$. Pour que $x$ appartienne à l'image essentielle du foncteur $\unsure{\varphi|\psi}$, il f. et s. que la sous-catégorie pleine $\underline S_{(\pi)/x}$ de $\deleted{\unsure{M_\pi}}\deleted{\unsure{M_\pi}}$ la catégorie $M_{\pi/x}$ (qui est filtrante et $=$ gde devant $\pi$) soit cofinale (donc aussi gde devant $\pi$).

b) Pour que $\underline S$ soit $\pi$-accessible dans $M$, il faut et il suffit que pour tt $x\in\mathrm{Ob}\,M$, on ait

\[x\in\mathrm{Ob}\,\underline S\quad\Longleftrightarrow\quad\underline S_{(\pi)/x}\ \text{est cofinale dans }M_{\pi/x},\]

\inserted{[ou encore, que $\underline S$ satisfasse à $L_\pi$, et que pour tt $x$ de $M$, on ait $x\in\mathrm{Ob}\,\underline S\Longrightarrow\underline S_{(\pi)/x}$ cofinal dans $M_{\pi/x}$]}

\uline{Démonstration}. La partie b) est une conséquence tautologique de a), que je vais prouver.

\deleted{\unsure{Si}}\inserted{Dire que} $x$ \deleted{\unsure{est}} est dans l'image essentielle de $\underline S$, signifie que \inserted{$(\unsure{\alpha|*})\ x\simeq\varinjlim_{i\in I}x_i$}, avec $I$ \unsure{grd|gd} devant $\pi$ et les $x_i$ dans $\underline S_{(\pi)}$. On sait que \deleted{\unsure{de}} \unsure{ceci|ces|ça} \unsure{implique|impliquent} ($I$ étant filtrant, et $M$ étant équivalent à $\mathrm{Ind}_\pi(M)$ puisque $\pi$ utile pour $M$) que $I\longrightarrow M_{\pi/x}$ est cofinale (4.4.2 page 79). Or ce foncteur se factorise par le foncteur d'inclusion pleinement fidèle $\underline S_{(\pi)/x}\unsure{\lhook\joinrel\relbar|\lhook\joinrel\longrightarrow} M_{\pi/x}$, lequel est cofinal \unsure{sss|ssi} il satisfait la condition $\mathrm{F}1$. Il est immédiat que cette condition est satisfaite, quand elle est satisfaite pour le foncteur composé en question, ce qui est bien le cas

% archive sheet number: 579
